% applications/contribution-DE.tex
\chapter{Beitrag}
\label{chap:applications-contribution}

Dieses Kapitel fasst die Konsequenzen des Beitrags der Arbeit aus
Anwendungssicht zusammen.

\section{Umfang des Beitrags}

Der zentrale Beitrag ist keine neue isolierte Anwendungsmethode, sondern
ein kohärentes, Alignment-zentriertes Ingenieurframework, das durch
kanonische Begriffe Entscheidungsunterstützung auf Anwendungsebene
ermöglicht.

\section{Beiträge auf Anwendungsebene}

Aus Anwendungssicht leistet die Arbeit folgende Beiträge:
\begin{itemize}
\item eine konsistente Entwicklung vom Begriff zur Anwendung,
\item operatorgestützte Laufzeit- und realisierungsbewusste Auswertung,
\item die ausdrückliche Trennung von Ingenieurobjekten und
Repräsentations- beziehungsweise Austauschformen
\item sowie praktische Interoperabilitätswege für BIM/GIS-orientierte
Arbeitsabläufe.
\end{itemize}

\section{Konsequenz für die Ingenieurpraxis}

Anwendungen werden reproduzierbar und vergleichbar, weil sie aus einem
stabilen Begriffskern statt aus projektspezifischer Ad-hoc-Semantik
abgeleitet werden.

\section{Verbindung zu AIM}
\begin{summary}
Der Beitrag des Anwendungsteils besteht darin zu zeigen, wie genehmigte
AIM-Begriffe zu konkreten Ingenieurfolgen führen, darunter
laufzeitgestützte Auswertung, realisierungsbewusster Vergleich und
interoperable Repräsentation zur Entscheidungsunterstützung.
\end{summary}
